\section{Software design} \label{sec:design}

This section describes the architectural decisions that turn twelve heterogeneous estimators into interchangeable components of a single workflow. The key idea is that every estimator returns the same result object and that every downstream tool consumes that result object rather than the estimator that produced it. The result is a clean separation between estimation and inference, which means that any improvement to the inference layer benefits every estimator at once.

\subsection{The Estimator protocol and the EstimationSummary contract}

Every estimator in the library implements the \code{Estimator} protocol defined in \code{econirl.estimation.base}. The protocol requires a \code{name} property and an \code{estimate} method with the signature
\begin{quote}
\code{estimate(panel, utility, problem, transitions) -> EstimationSummary}.
\end{quote}
The four arguments specify the data, the utility function, the structural primitives, and the transition tensor. The return value is a single object that carries the parameter point estimates, the parameter names, the standard errors, the variance-covariance matrix, the Hessian at the optimum, the converged value function, the choice probabilities, the goodness-of-fit measures, and the identification diagnostics. The same object exposes formatted output through a \code{summary} method and offers \code{wald\_test}, \code{confidence\_interval}, \code{t\_statistics}, \code{p\_values}, \code{to\_latex}, and \code{to\_dataframe} as convenience methods.

A subclass of \code{BaseEstimator} implements only the optimization core via the abstract \code{\_optimize} method, which returns an \code{EstimationResult} with the raw fit. The base class then computes standard errors, identification diagnostics, and goodness-of-fit, and assembles the final \code{EstimationSummary}. Adding a new estimator requires writing the optimization loop and nothing else.

\subsection{The shared inference layer}

The inference layer lives in \code{econirl.inference} and is shared across all estimators. Four standard error methods are available through the single entry point \code{compute\_standard\_errors}. The asymptotic method inverts the negative Hessian. The robust sandwich method combines the inverse Hessian with the outer product of the per-observation score contributions. The nonparametric bootstrap resamples individuals with replacement, refits the estimator on each resample, and reports the empirical standard deviation across replications. The clustered method sums the gradient contributions within each individual before forming the outer product, which is the appropriate adjustment for panel data with within-individual dependence. The user selects a method by passing a single string at estimator construction time. All four methods return the same \code{StandardErrorResult} structure.

The identification layer in \code{econirl.inference.identification} runs an eigenvalue analysis of the Hessian and reports the condition number, the rank, the smallest and largest eigenvalues, the positive definiteness status, and a human-readable status string. A separate diagnostic function \code{diagnose\_identification\_issues} flags problematic parameters by checking for near-zero eigenvalues and high pairwise correlations exceeding 0.95. These diagnostics are computed automatically as part of every \code{estimate} call and are stored in the \code{identification} field of the returned \code{EstimationSummary}.

The profile likelihood module in \code{econirl.inference.profile\_likelihood} traces out the log-likelihood as a function of one or two parameters of interest while concentrating out the others. The hypothesis test module in \code{econirl.inference.hypothesis\_tests} provides Wald tests, likelihood ratio tests, and Hausman tests through a uniform interface that accepts any \code{EstimationSummary}. The bootstrap module in \code{econirl.inference.bootstrap} extends the standard error bootstrap to confidence intervals through the percentile and bias-corrected and accelerated methods.

\subsection{Counterfactuals and simulation}

Counterfactual analysis consumes the fitted result rather than the estimator. The functions \code{counterfactual\_policy} and \code{counterfactual\_transitions} in \code{econirl.simulation.counterfactual} take an \code{EstimationSummary} and either a new parameter vector or a new transition tensor and return the implied choice probabilities and value function. The \code{elasticity\_analysis} function reports the partial derivative of any scalar policy summary with respect to any parameter through automatic differentiation. The \code{compute\_stationary\_distribution} function returns the long-run state distribution under the estimated policy, which is the workhorse object for welfare comparisons.

The synthetic data layer in \code{econirl.simulation.synthetic} provides \code{simulate\_panel} for forward simulation under a given policy and \code{run\_monte\_carlo} for parameter recovery experiments under a known data generating process. The Monte Carlo function returns the bias, the root mean squared error, and the empirical coverage rate of the standard errors across replications, which is the standard validation we run for every new estimator before release.

A small visualization layer in \code{econirl.visualization} renders policies, value functions, and counterfactual overlays from any \code{EstimationSummary} via \pkg{Matplotlib}.

\subsection{The JAX backbone}

The library is built on \pkg{JAX} \citep{jax2018github}. Automatic differentiation produces gradients and Hessians without hand-coded derivatives, which removes the most common source of bugs in structural code. Just-in-time compilation accelerates inner loops by an order of magnitude relative to \proglang{NumPy}, and vectorized mapping enables bootstrap resampling and Monte Carlo replications without manual thread management. Default precision is 64-bit for stable Hessian computation; the user can drop to 32-bit on a graphics processing unit, with the dispatch handled automatically.

